\documentclass[11pt,a4paper]{article}

\usepackage[LGR,T1]{fontenc}
\usepackage[utf8]{inputenc}
% Transliteraciones con macron+agudo (megalḗtōr, kynṓpidos...)
\DeclareUnicodeCharacter{1E17}{\'{\=e}}
\DeclareUnicodeCharacter{1E53}{\'{\=o}}
\usepackage[greek.ancient,spanish]{babel}
\usepackage[a4paper,margin=2.8cm]{geometry}
\usepackage{microtype}
\usepackage[hidelinks]{hyperref}

% Griego politónico escrito directamente en UTF-8:
% \gr{...} para palabras sueltas; bloques con \foreignlanguage{greek}{...}
\newcommand{\gr}[1]{\foreignlanguage{greek}{#1}}

\hypersetup{
  pdftitle={Leer la Odisea en tiempos iletrados. Octava parte: Circe}
}

\setlength{\parindent}{1.25em}
\setlength{\parskip}{0.35em}
\raggedbottom

\title{Leer la \emph{Odisea} en tiempos iletrados \\ Octava parte \\ Circe}
\author{}
\date{}

\begin{document}

\maketitle

¿Y si la gran pasión de Odiseo no hubiera sido Penélope, sino Circe? ¿Y si su viaje a Ítaca no hubiera sido un regreso al hogar, sino la fuga de una cárcel de amor?

Desde luego, no es eso lo que nos cuenta Nolan. Su episodio da muchos detalles bien observados, entre ellos la isla llena de falsos animales, o la transformación de los compañeros de Ulises en cerdos cuando aún conservan forma humana, reflejada por el ansia inhumana con la que comen. Pero su retrato de Circe sigue el mismo patrón que en el caso de Polifemo: prefiere presentarnos un personaje simple, casi plano, cuyas motivaciones son una mezcla de miedo y odio genérico a los hombres. A cambio de la superficialidad psicológica, nuestro fílmico director nos obsequia con toda una exhibición de ingeniería plástica, en la que la bruja transforma, \emph{a mano} y sudando lo suyo, a los hombres en cerdos.

Nada más. La relación entre Circe y Ulises es una simple transacción comercial ---mis hombres por tu hermana transformada en urraca---, adornada con unas cuantas frases de compromiso. Nolan no nos explica quién es Circe ---a pesar de que se trata del personaje más sobrenatural del relato, en cuanto que su magia es explícita, a diferencia de la de Calipso, por ejemplo---, ni nos informa de su agenda. La visita a la isla es breve y el viaje a los infiernos que sigue se muestra como un episodio independiente.

Veamos qué nos cuenta Homero.

Circe aparece en el canto X, pero no directamente. Antes, los aqueos llegan a Eolia, isla donde habita el dios de los vientos, junto a sus hijos e hijas ---que no solo están juntos sino incestuosamente revueltos---, todos ellos pasándoselo estupendamente, como ya sabemos que es costumbre de los inmortales:

\begin{verse}
Hasta que al fin la isla Eolia tocamos; en ella vivía\\
Eolo Hipótada, amado del cielo. Era aquella una isla\\
flotante y toda a redor por un muro de bronce ceñida\\
inquebrantable, y alzadas de roca que erguíanse lisas.\\
En su palacio doce hijos nacieron, y allí los tenía;\\
seis hijas tuvo, y seis hijos, y, entrados en flor de la vida,\\
a cada hijo le vino a dar él por esposa una hija.\\
Junto a su padre querido y su honrada madre, los días\\
ellos en fiesta los pasan, que siempre hay mil viandas servidas\\
y las paredes rezuman contento y buen husmo a cocina;\\
eso de día, que luego la noche a cada uno lo arrima\\
al lecho junto a su santa con los cobertores por cima.
\end{verse}

Eolo ofrece hospitalidad y los navegantes aprovechan para tomarse un respiro:

\begin{verse}
A su ciudad, pues, llegamos, y a aquellas estancias magníficas.\\
Un mes entero me tuvo él allí, que de todo quería\\
detalle: Ilión, el regreso de aqueos, las naves argivas…\\
Yo de eso todo, y a ley, por menudo el relato le hacía.
\end{verse}

Cuando Odiseo le pide permiso para continuar su viaje, el simpático dios no solo se lo da, sino que le ofrece un regalo, como ya sabemos que exige la \emph{xenía}, cumplida aquí al pie de la letra. Eolo le regala a Ulises nada menos que la bolsa de los vientos. Es un tesoro que le permitirá llegar a su patria en un santiamén, pero, curiosamente, su magia exige no dejar que nadie se acerque a ella y, desde luego, no abrirla ---estamos, claro, ante otra versión del mito de Pandora---. Lo que ha hecho el dios del viento es aprisionar a todos sus díscolos súbditos en un saco muy bien cerrado, para que no molesten durante la travesía.

\begin{verse}
Cuando le hablé de seguir ya camino y rogué despedida,\\
él no me dijo que no, y hasta me preparó la partida.\\
Tras desollar un buey nueveañero y hacerme una odrina,\\
tráfago en ella de vientos hinchados me ató que mugían,\\
pues despensero de vientos lo había hecho a él el Cronida\\
y en él estaba excitar o calmar al ciclón que quería.\\
Luego en la nave con cuerda de plata otros nudos le hacía\\
porque ni por las costuras soplara pizca de brisa.\\
Nos levantó un sostenido poniente al salir de la isla\\
para llevar a navíos y hombres, mas no llegaría\\
eso a cumplirse: nuestro desatino forjó nuestra ruina.
\end{verse}

¿Qué insensatez van a cometer Ulises y los suyos esta vez? Una de juzgado de guardia:

\begin{verse}
Hicimos rumbo muy bien, noche y día, por nueve jornadas\\
cuando en la décima se apareció como luz nuestra patria,\\
y hasta veíamos los que atendían allá sus fogatas.\\
Y de repente ahí se entró el dulce sueño en mis fuerzas cansadas,\\
que es que fui siempre a la escota del barco, y no quise dejarla\\
a hombre ninguno, por antes mi tierra poder alcanzarla.
\end{verse}

Ya está, se dice Odiseo: ya han llegado, todo ha salido bien y nada le impide echarse una siesta para estar despejado cuando arriben al puerto. Mala idea:

\begin{verse}
Mis compañeros, entonces, cambiaban entre ellos palabras,\\
y se decían que yo me llevaba oro y plata a mi casa\\
que del magnánimo Hipótada Eolo, seguro, eran dádiva.\\
Y así uno entre ellos decía cruzando con otro miradas:

«Ay, hay que ver cómo quieren y cómo honran todos a este,\\
sea cual sea la tierra y ciudad a la que se acerque.\\
Muchos tesoros de Troya se trae en botín, y excelentes;\\
y ya lo ves, hecho el mismo viaje, a nosotros, su gente,\\
con nuestras manos vacías a casa llegando nos tiene.\\
Y encima ahora se lleva en amor y compaña el presente\\
de Eolo. Pronto, vayamos a ver lo que sea que encierre,\\
la tanta plata y el oro que el cuero, seguro, guarece».
\end{verse}

¡Vaya encanto de tripulación! Nolan captura bastante bien en su film la naturaleza levantisca de los compañeros del héroe, pero Homero es mucho más despiadado: los pinta ingratos, envidiosos e imprudentes. Y por supuesto, se buscan su propia ruina:

\begin{verse}
Así decían. Venció el mal consejo de mis compañeros,\\
y desataron el odre, y aupáronse todos los vientos.\\
Un huracán que sobre ellos cayó los llevó mar adentro\\
llorando, porque la patria se les marchaba. Yo, en esto,\\
me desperté y en mi pecho sin tacha me vi discurriendo\\
si de la nave saltar y acabar de una vez en el piélago,\\
o soportarlo en silencio y seguir con los vivos más tiempo.\\
Y soporté, y me quedé, y me tendí, mi cabeza cubriendo\\
sobre las tablas, mientras arrastraba el ciclón de mal viento\\
a la isla Eolia las naves de nuevo, y mis hombres gimiendo.
\end{verse}

Los versos son estremecedores. El engañador engañado, el traidor traicionado por su propia gente. Y de paso, otro mito deslizándose entre los versos, el de Sísifo, que empuja su piedra montaña arriba solo para verla rodar por la ladera cuando ya casi había alcanzado la cumbre, igual que Ulises ve alejarse Ítaca cuando ya se imaginaba fondeando en su puerto.

De vuelta a Eolia, Odiseo trata de suplicar clemencia, humillándose ante Eolo:

\begin{verse}
Desembarcamos allí para hacer de agua dulce el abasto,\\
y por las lévedas naves comiendo enseguida ya estábamos.\\
En cuanto todos de yanta y bebida la gana amansamos,\\
a un compañero conmigo tomé y a otro más como heraldo\\
y fui al glorioso palacio de Eolo; a él allí festejando\\
me lo encontré con su esposa y su prole. No entré en el palacio,\\
que, sin cruzar el umbral, ante el mismo portón nos sentamos.\\
Y se asombraron aquellos al vernos y nos preguntaron:

«¿Cómo, Odiseo, tú aquí? ¿Qué demonio se te ha echado encima?\\
Te despedimos de gana y de fe porque así llegarías\\
a esa tu patria y tu casa, o adonde quisieras venida».

Eso dijeron, y yo les hablé desde mi desconsuelo:\\
«Daño me dio la mi mala compaña, y daño el mal sueño;\\
mas remediádmelo, amigos, que está en vosotros hacerlo».
\end{verse}

Pero esta vez no hay suerte. Eolo considera que su antiguo huésped está maldito y lo echa a patadas:

\begin{verse}
«Sal de la isla ya mismo, que vivo no lo hay más maléfico;\\
yo permitido no tengo amparar ni mandar de regreso\\
a hombre que sea de dioses benditos aborrecimiento.\\
Sal, que si has vuelto tú aquí, como odiado del cielo lo has hecho».
\end{verse}

Y los aqueos, con el corazón «sombrirroto», salen a remo, ya que no sopla una brizna de viento:

\begin{verse}
Así diciendo, me echó de su casa. Y yo en hondo sollozo.\\
De allí otra vez navegamos con el corazón sombrirroto,\\
mientras roía a los hombres su espíritu el remo penoso\\
por necedad, que ya allí no asomaba la escolta de un soplo.
\end{verse}

Por necedad, todo por necedad. A veces da la impresión de que la estupidez humana es el tema central de la \emph{Odisea}.

Al cabo de seis días, los navegantes llegan a la isla de los lestrigones, que resultan ser unos energúmenos que casi se los comen crudos. Nolan refleja el episodio haciendo aparecer unos gigantones vestidos con armadura plateada que primero te cortan la cabeza y después te dan los buenos días, pero en realidad se ciñe bastante a Homero, que también los pinta como montañas ambulantes ---esencialmente una versión armada y bien organizada de los cíclopes--- y, en todo caso, no hay mucho desarrollo psicológico en el pasaje. Los lestrigones son agentes de destrucción que diezman la tripulación de Odiseo y le hunden todos los barcos menos el suyo.

\begin{verse}
Mientras adentro del puerto honduroso hacían ellos estrago,\\
tirando yo de la espada buída que va a mi costado,\\
cable enseguida con ella corté del navío proiblavo.\\
Y presuroso a mis hombres moví y ordené sin descanso\\
jalar de remos para sortear ese tanto quebranto;\\
y sacudían el mar todos ellos, de muerte espantados.\\
A la maralta escapó felizmente de aquellos peñascos\\
mi nave; todas las otras allí su final encontraron.
\end{verse}

Diezmados y exhaustos, tras haber escapado por los pelos, llegan a Eea, la morada de Circe:

\begin{verse}
Y así llegamos a la isla de Eea; en ella habitaba\\
Circe beltrenza, la fascinadora deidad vocihumana,\\
que era de Eetes mientesfunestas verísima hermana;\\
los dos nacidos del Sol, que la luz a mortales regala,\\
y ambos por madre venidos de Perse, de Océano alhaja.
\end{verse}

Fíjate, observadora lectora, detallista lector, cómo nos presenta Homero a la bruja beltrenza. La palabra clave que la define es \gr{δεινή} (deinḗ), literalmente temible o terrible, pero también imponente y venerable, términos que Tobal condensa en «fascinadora». También es importante «vocihumana», traducción afortunada de \gr{αὐδήεσσα} (audḗessa) que asegura que Circe es capaz de hablar con los hombres (y de embaucarlos, como veremos) en su propia lengua. Y, por supuesto, en jerarquía, la terrible deidad de las bellas trenzas está por encima de Calipso y casi de cualquiera. Hija del Sol y de la deidad marina Perse (o Perseis en Hesíodo). Es decir, una diosa de pleno derecho.

Recién llegados, Ulises se va a explorar, distingue el palacio de Circe ---nada de cabaña--- en lo alto de una peña y luego caza un ciervo con el que él y sus hombres se dan un festín en la playa. Al día siguiente les anima a que vayan a explorar, pero la tripulación, después de habérselas visto con cíclopes y lestrigones, no está para muchas aventuras. Finalmente forman dos grupos. Uno al mando del propio Odiseo y el otro capitaneado por su lugarteniente Euríloco. Echan a suertes a quién le toca explorar y la suerte cae en este último, que parte con muy pocas ganas:

\begin{verse}
Conque marchó, y juntamente con él veintidós compañeros\\
llorando; atrás nos dejaban allí a nosotros gimiendo.
\end{verse}

Viene ahora la descripción de la fauna que rodea el palacio:

\begin{verse}
Dieron por una honda hoz con la casa de Circe, unos lienzos\\
hechos de piedra pulida, y, en torno, un lugar bien abierto.\\
Y por allí dormitaban leones y lobos serreños\\
que ella tenía hechizados por obra de algún mal veneno.\\
Pues, sin amago de echarse a mis hombres, en pie se pusieron\\
para halagarlos meneando sus colas al paso de ellos.\\
Y como, cuando al volver de una cena su amo, los perros\\
mueven el rabo, pues siempre él les trae golosina, contentos,\\
así los leones allí y los lobos pezuñirrecios\\
fiestas hicieron; miraban aquellos los mostros con miedo.
\end{verse}

Nolan nos muestra los animales, sugiriendo que se trata de humanos convertidos, como también lo es el ciervo que caza Ulises. Homero no lo deja claro ---desde luego el ciervo no se convierte en humano---, pero el detalle no es importante. Las pacíficas fieras revelan que los aqueos se están acercando a la morada de una hechicera peligrosa. Pero los siguientes versos dan un giro al argumento.

\begin{verse}
Ya en el portón de la trencigalana deidad, se tuvieron,\\
y a Circe oían cantar ---era hermosa su voz--- allí dentro,\\
puesta al telar a una tela inmortal ---de una diosa era empleo---\\
fina y graciosa y brillante, que así son las obras del cielo.
\end{verse}

No hay en Homero una granjera que vuelva a su choza acarreando sus bártulos, sino una deidad de hermosa voz ---el detalle me impactó tanto que mi propia versión de Circe se construye a partir de esa voz--- tejiendo una tela inmortal.

¿Suena familiar? Por supuesto. Penélope y Circe son tejedoras. La simetría con que Homero nos presenta a estas dos formidables mujeres no puede estar más clara. Pero también hay una asimetría, como escribe Candelas Gala en su epílogo a \emph{Abandonando Ítaca}\footnote{\url{https://eolasediciones.es/libro/abandonando-itaca-leaving-ithaca/}}:

\begin{quotation}
Sí, porque Penélope y Circe pueden verse como dos caras de la misma moneda. Ambas son tejedoras, pero de un tipo muy diferente. Si damos crédito a los relatos clásicos de Circe en la cultura occidental, se nos aconsejaría alejarnos de una mujer peligrosa, malvada, el arquetipo de una hembra depredadora. Su representación como bruja ha cambiado un poco con el tiempo, cuando se cuestionó la existencia de las brujas, pero ese cambio no la ha librado de una reputación negativa como mujer delirante y sexualmente libre. Cambió a los hombres de Ulises por cerdos, como si se hubiera anticipado al apelativo moderno de «cerdos» que suelen utilizar las mujeres que reaccionan contra el excesivo machismo patriarcal. Pero también dio a Ulises consejos útiles sobre el descenso al Hades para adquirir conocimientos sobre los caminos de los dioses y cómo volver a casa sano y salvo. Ella misma era una gran conocedora de las hierbas y las pociones, una \emph{polifarmakos} según uno de sus epítetos homéricos. Su nombre evoca el círculo de Circe, que representa en su rueca cósmica. Homero la llamó Circe de las trenzas, sugiriendo en los nudos y trenzas de su cabello las fuerzas de creación y destrucción. Su tejido tenía el poder de gobernar las estrellas y los destinos de los hombres.

Penélope también era hilandera, enhebrando un sudario que tejía y destejía repetidamente. Mientras que el hilado de Circe era transformador y cambiante, el círculo vicioso de Penélope de repetición de lo mismo no podía competir con el tejido imprevisible y siempre renovable de Circe.
\end{quotation}

Pero no nos adelantemos a Homero. Uno de los aqueos, Polites, a quien Odiseo llama «el más atento y leal de los compañeros», propone que hagan ruido para manifestar su presencia. Circe los oye y los invita a pasar:

\begin{verse}
Ella enseguida salió y, abriendo sus puertas brillantes,\\
los invitaba a pasar; iban todos tras ella, ignorantes…\\
Tan solo Euríloco fuera quedó, presintiendo un mal trance.\\
Ellos adentro, los hizo por tronos sentarse y sitiales,\\
y con cebada, con queso y miel verde aliñaba un brebaje en\\
vino de Pramnos, en el que ligó para el mismo potaje\\
drogas siniestras que hacían olvido del lar de los padres.\\
De que les dio la poción y bebieron, golpeó en ese instante\\
con su varita, y a las cochiqueras los lleva a guardarse.\\
Que de cochinos tenían la testa, la voz, la pelambre\\
y la figura, mas eran sus mientes las mismas de antes.
\end{verse}

Fíjate en el detalle, biensabida lectora, siempreatento lector. Homero despacha la transformación en cerdos en unos pocos versos. Circe recurre a «drogas siniestras que hacían olvido del lar de los padres» y aquí encontramos otro de los interminables ecos que resuenan en la \emph{Odisea}\footnote{Como una moneda que va cayendo al fondo de un pozo, decía Pessoa.}. El filtro de Circe hace olvidar y evoca el de Helena, que impide sentir tristeza. Estamos ante poderosas hechiceras armadas con filtros ---y varitas mágicas---. ¿Hacía falta la elaborada escena de ingeniería plástica en la que la bruja transforma a los hombres de Ulises en gorrinos, como quien dice, «a mano»? ¿No podía el bienquerido director haber dedicado todos esos minutos a darle más fondo al personaje? Otro detalle importante, que pone de manifiesto la crueldad de Circe, es que los cerdos lo son solo en apariencia, mientras que sus mentes permanecen intactas.

En la \emph{Odisea}, Euríloco no es transformado en cerdo, sino que escapa y le da noticias a su capitán.

\begin{verse}
«Como ordenaste, cruzamos el bosque, preclaro Odiseo;\\
dimos por una honda hoz con casa hermosa, de lienzos\\
hechos de piedra pulida, y, en torno, un lugar bien abierto.\\
Alguien labraba al telar una tela cantando allí dentro,\\
una mujer o una diosa. A voces llamaron aquellos.\\
Ella enseguida salió, les abrió los portones espléndidos,\\
y los invita a pasar; todos iban detrás, bien ajenos…\\
Yo quedé fuera, que pálpito tuve de trance no bueno.\\
Todos, a un tiempo, invisibles se hicieron, y ni uno de ellos\\
apareció; que yo allí mucho rato me estuve en acecho».
\end{verse}

Nada de chozas, nada de mujeres desconfiadas. Una casa hermosa de piedra pulida, una diosa aún más hermosa que canta. ¿Por qué insiste tanto Homero en la música de Circe? ¿Qué quiere decirnos? Ulises se dispone a explorar y le pide a su lugarteniente que le acompañe, pero este, muerto de miedo, le pide quedarse. Nuestro héroe, al que nadie puede negarle el arrojo, se dirige a casa de la bruja, cuando el dios Hermes se le cruza en el camino:

\begin{verse}
Mas cuando, haciendo camino por la hoya encantada, me hallaba\\
ya de esa Circe milpocimera llegando a la casa,\\
Hermes varitadeoro, saliéndome al paso, me asalta\\
---la casa ya ante mis ojos--- igual que un muchacho en estampa\\
recién de barbiponiente, en la flor de la edad más graciada;\\
y me tomó de la mano, y así por mi nombre me hablaba:
\end{verse}

Hermes no solo le pone sobre aviso, sino que le da los medios de enfrentarse a la hechicera:

\begin{verse}
«¿Dónde, cuitado, otra vez marchas solo por estos rincones,\\
y sin saber lo que hay? Por la casa de Circe tus hombres\\
en cochiqueras cerradas guardados están, son cebones.

¿Vienes aquí a liberarlos? Te digo lo que aún no conoces:\\
téntalo y no volverás, quedarás donde están esos pobres.\\
Ah, pero vengo a salvarte y librarte de sus maldiciones.\\
Sí, vete a casa de Circe, mas este remedio antes tómate\\
benigno, que alejará de tu testa la hora sin nombre.

Voy a contarte el hacer pernicioso de Circe, y por su orden.\\
Preparará una mixtura y pondrá en la comida ese cóctel;\\
mas ni con eso hechizarte podrá, que a ese pisto se impone\\
la droga buena que ahora te doy; luego harás lo que oyes.

Cuando con larga varita intente ella darte un azote,\\
saca la espada aguzada que llevas al muslo y, de golpe,\\
a perseguirla, como deseando matarla, te pones.

Ella, temblando, a su cama invitarte querrá: nunca nones\\
debes a lecho de diosa decir si de cierto a tus hombres\\
los quieres ver liberados y que ella con mimo te aloje;\\
mas que te jure con gran juramento bendito de dioses\\
que ella no va a maquinar contra ti mal ninguno a mayores,\\
no se te lleve la hombría y valor cuando te desarropes».
\end{verse}

Observa, detallista lectora, minucioso lector, que la aparición de Hermes no es casual. Odiseo es un simple mortal y es canónico que un mortal no puede enfrentarse solo a una diosa, lo que hace imprescindible que otro olímpico equilibre las tornas. El texto es delicioso. Frente al veneno, antídoto. Frente a la varita, espada. ¿Y qué hará la Circe cuando se vea acorralada? ¡Intentará encamar a Odiseo, por supuesto! Lo mejor de todo es que Hermes no le impone al héroe que se niegue, sino que le advierte de que antes de meterse en el tálamo con la divina, se asegure de que le restituya a sus hombres y le garantice hospitalidad. Nadie aquí parece preocupado por el hecho de que el valiente guerrero eche una cana al aire mientras Penélope espera, tejiendo y destejiendo su mortaja.

Odiseo llega a casa de Circe y empieza la acción:

\begin{verse}
Ya en el portal de la trencigalana deidad, me detengo,\\
y allí parado le grito, y la diosa mi voz ya está oyendo.\\
Y ya sus puertas brillantes me abría saliendo en un vuelo\\
y me invitaba a pasar; la seguí con el ánimo inquieto.\\
Dentro, me lleva a sentar a un trono argenticlaveño\\
bello, de buena labor, con un lindo alzapié sobre el suelo.\\
Y me aliñó ella un brebaje en un cáliz de oro poniendo\\
en él la droga, tejiendo mi mal en su pecho malevo.\\
Y me lo dio, y lo bebí, pero no, no hubo allí encantamiento\\
por más que con su varita me diera, este ensalmo diciendo:

«Ve a la pocilga ya mismo, y acuéstate con tu pandilla».\\
Dijo, y, sacando yo entonces de al muslo mi espada buída,\\
me abalancé sobre Circe como por quitarle la vida.\\
Ella, gritando, se me agachó y se abrazó a mis rodillas\\
y entre sollozos palabras aladas así me decía:
\end{verse}

\begin{verse}
«¿Quién y de qué gentes tú? ¿Y tus padres? ¿En qué ciudad vives?\\
No entiendo cómo, aun bebida la pócima, embrujo no hice;\\
jamás un hombre ---jamás--- a esta droga logró resistirse\\
cuando, al beber, traspasó ella el redil de los dientes. Bebiste,\\
pero propósito hay dentro de ti que el hechizo no rinde.\\
Tú eres aquel milerrante Odiseo, el que siempre me dice\\
el Argifontes varitadeoro que habrá de venirme\\
un día aquí desde Troya en su léveda nave; y viniste.\\
Guarda esa espada en la vaina, y subamos tú y yo sin más quite\\
a nuestra cama, porque se nos venga que del confundirse en\\
lecho y amor tú confíes en mí y en ti yo confíe».
\end{verse}

Tal y como había predicho Hermes. Ulises, que juega con ventaja, no se fía y le pide un juramento:

\begin{verse}
Eso me dijo, y entonces hablé para así responderle:\\
«¿Cómo me puedes pedir que a ti, Circe, yo amable me acerque\\
si de mis hombres has hecho cochinos entre estas paredes\\
y a mí me tienes aquí y con artera cabeza me ofreces\\
entrar en tu dormitorio para que contigo me acueste y\\
puedas, cuando me desnude, la hombría quitarme y el temple?\\
No, no querría a tu cama subir, por más que te pese,\\
salvo que pases tú, diosa, por un juramento solemne\\
de que alumbrando no estás contra mí nuevo mal sobre este».
\end{verse}

A la hechicera le falta tiempo para jurar, a Ulises para perdonarla y a Homero para meterlos en la cama:

\begin{verse}
Dije, y juró ella enseguida, como se lo estaba pidiendo.\\
Y cuando ya hubo jurado y cerrado quedó el juramento,\\
con Circe entonces subí a su bienhermosísimo lecho.
\end{verse}

¡Qué duros los trabajos del héroe! Mientras la comodona de su esposa se entretiene tejiendo y destejiendo, a él no le queda otra que darse unos revolcones divinos en el bienhermosísimo lecho de Circe.

Homero nos cuenta cómo las criadas de Circe lavan a Odiseo, lo perfuman, lo ungen con aceite y lo sientan a la mesa. Ulises, por una vez, parece desganado. La diosa le dice entonces:

\begin{verse}
«¿A qué te sientas así, Odiseo, como un sin palabras\\
y pan y vino no tocas, mientras te recomes el alma?\\
¿Piensas que habrá un otro engaño? Pues de eso no debes en nada\\
desconfiar: me di a un juramento que es jura sagrada».
\end{verse}

Lo cierto es que el héroe hace bien en no fiarse de los juramentos de los inmortales, tan aficionados a hacer promesas como a saltárselas. Además, Hermes le ha advertido que exija el regreso de sus hombres como prueba de buena voluntad.

\begin{verse}
Dijo ella así, y respondiéndole yo le decía de vuelta:\\
«¿Ay, Circe, qué hombre iba a haber ---y digo hombre cabal--- que tuviera\\
de atiborrarse el valor, entregado a comienda y bebienda,\\
sin que libere a sus hombres primero y sus ojos los vean?\\
Que si me pides de corazón que yo coma y que beba,\\
suéltalos, y que a mis ojos mis buenos leales me vuelvan».
\end{verse}

La hechicera no se hace de rogar, transforma a los compañeros de nuevo en hombres e invita a Ulises y a los suyos a arrastrar las naves hasta el arsenal de la isla y a instalarse con ella. Ulises la obedece, regresa a la playa y le dice a su tripulación:

\begin{verse}
«Al arenal arrastremos la nave primero y, ya en firme,\\
a los tesoros y jarcias daremos en cueva escondite;\\
luego aprestaos vosotros ---y no falte nadie--- a seguirme\\
porque a los vuestros veáis en la casa divina de Circe\\
bebiendo a bien y comiendo, que hay más de lo que un año rinde».
\end{verse}

Euríloco, como era de esperar, no se fía y dice así:

\begin{verse}
«Ay desgraciados, ¿ir dónde? ¿Por qué de más mal en amores?\\
¿Ir a la casa de Circe, que a todos nosotros a un golpe\\
en porcachones o en lobos nos convertirá o en leones,\\
para a la fuerza allí ser de esa su alta mansión guardadores?\\
Ved lo del Cíclope, cuando a su aprisco se entraron los hombres\\
---vuestros iguales--- con el temerario Odiseo a mayores:\\
la altanería de este de aquellos movió la hecatombe».
\end{verse}

Tiene toda la razón, obviamente. ¿Y cuál es la reacción de Ulises?

\begin{verse}
Dijo, y por mi pensamiento pasó, en esas dudando,\\
tirar del sable hojiluengo que va al recio muslo ajustado\\
y, tras cortar su cabeza, echarla a la tierra rodando,\\
y eso que me era pariente cercano; pero me frenaron\\
uno tras otro mis compañeros con estos halagos:

«Hijo del cielo, dejemos a este, si tú lo permites,\\
y aquí se quede junto a la nave y que nos la vigile;\\
y a los demás llévanos a esa santa morada de Circe».

Esto diciendo, dejaban el mar, el navío dejaban.\\
Pero ni Euríloco quiso quedarse en la nao concavada,\\
y nos seguía; temiendo a bien cierto mi atroz amenaza.
\end{verse}

¡Un encanto este hombre! Euríloco es su lugarteniente, su amigo del alma y se ha librado por un pelo, nunca mejor dicho, de que le cortara la cabeza por atreverse a llevarle la contraria. El caso es que vuelven a palacio y Circe les reitera a todos la invitación a pasarlo bien y descansar después de tantas cuitas. Los aqueos aceptan y se quedan... ¡por un año! Claramente, no lo pasarían muy mal. Pero al cabo de tantos meses, les empieza a remorder la conciencia y le dicen al capitán:

\begin{verse}
«Hombre de dios, del solar de la patria tendrás que acordarte\\
ya si la voz de los cielos dictado dejó que te salves\\
y que a tu casa techalta al fin llegues y al lar de los padres».
\end{verse}

Ulises, por fin, recuerda que lo suyo es volver a casa y dejarse de juergas.

\begin{verse}
Yo subí entonces de nuevo a la cama relinda de Circe\\
y a sus rodillas le supliqué, y la diosa a mí oíame,\\
y, levantando de verbos alada la voz, por fin dije:

«Cúmpleme, Circe, ahora eso de lo que me hicieras promesa,\\
mandarme a casa; que el alma se me anda por eso ya inquieta,\\
y la de mis compañeros, que mi corazón atormentan\\
aquí y allá lamentándose en cuanto que tú no estás cerca».
\end{verse}

La diosa parece consentir, pero pone una condición tremenda:

\begin{verse}
«Sangre Laercia, Odiseo milmañas, hijo del cielo,\\
no os quedéis ya por más tiempo en mi casa si ya no hay deseo.\\
Mas fuerza es que otro viaje cumpláis y destino primero\\
adonde tienen la horrenda Perséfone y Hades su reino\\
por del aliento tener del tebano Tiresias agüero,\\
aún con el seso en su sitio, que solo a este ciego agorero,\\
cuando murió, le mantuvo Perséfone su entendimiento\\
y lucidez, en un mundo de sombras en revoloteo».
\end{verse}

Circe manda a su amante a hablar con Tiresias en el Hades, ese sitio horrible donde iban las almas de los mortales para convertirse en sombras sin entendimiento. Los versos son espeluznantes: «solo a este ciego agorero, / cuando murió, le mantuvo Perséfone su entendimiento / y lucidez, en un mundo de sombras en revoloteo».

A Ulises, claro, la propuesta no le hace ninguna gracia:

\begin{verse}
Eso me dijo, y a mí el corazón se me desmoronaba.\\
Y allí sentado en el lecho me puse a llorar y mi alma\\
ya no quería vivir ni quería del sol ver más llama.\\
Mas cuando de revolverme y llorar se colmaron las ganas,\\
tuve palabras entonces para contestar su demanda:

«Circe, ¿y quién va a conducirnos ahora por ese camino?\\
Nadie ha llegado a la casa de Hades en negro navío».
\end{verse}

La hechicera le explica que será un viento mágico quien le lleve.

\begin{verse}
Hijo del cielo, Odiseo milmañas, sangre Laercia,\\
pierde cuidado, por falta de guía en la nave no hay pena;\\
fíjale el palo sin más y que blancas le caigan las velas\\
y espera solo: verás cómo el soplo del Bóreas la leva.\\
Solo después de cruzar el Océano con tu goleta,\\
en punta baja hallarás los bosques de Persefonea,\\
álamos altos y sauces machíos entre cañaveras,\\
junto al Océano remolinhondo allí el barco fondea y\\
de la mansión aguanosa de Hades darás en la senda.
\end{verse}

Le explica también lo que tiene que hacer al llegar:

\begin{verse}
Al Aqueronte el Piriflegetonte sus rápidos echa\\
con el Cocito, que es brazo de Estigia la malagunera,\\
y hay una peña donde ambos dos ríos soneros se encuentran.

A ella tú acércate, héroe, y ---escúchame y ripio no pierdas---\\
abre de un codo a lo largo y de un codo a lo hondo una cuenca\\
para que viertas por todos los muertos en ella una ofrenda:\\
primero de lechimel, dulce vino al seguido, y, sobre ellas,\\
una tercera de agua; y harina blanca le avientas.

Invocarás de los muertos, de hinojos, las nulas cabezas,\\
y les dirás que ya en Ítaca harás de una vaca horra y buena\\
un sacrificio en tu casa y el lar llenarás de preseas,\\
y que un carnero, además, para él solo y aparte, a Tiresias\\
le inmolarás, todo negro, de vuestro rebaño excelencia.

Cuando a las ínclitas razas de muertos rezadas las tengas,\\
degollarás un cordero tú allí y una oveja bien negra\\
con la testuz inclinada al Erebo; tú vuelve tu testa\\
como a seguir las corrientes del río: allí ya la cuenta\\
vas a perder de las sombras de muertos difuntos que llegan.

En ese mismo momento a tus hombres apremia y ordena\\
que, desollando las reses que a crudo bronce en la tierra\\
tú degollaste, las quemen, rezándole al dios y a la dea,\\
a Hades tremendo y con él a la hórrida Persefonea.

Y tú, de la hoja tirando que llevas al muslo sujeta,\\
mira que muerto ninguno acerque su huera cabeza\\
a aquella sangre hasta que de Tiresias te venga respuesta.\\
Presto, caudillo de hombres, allí acudirá ese profeta\\
que te dirá el buen camino y el tiempo de tu carrera\\
y ese regreso que por la llanura de peces tú tengas.
\end{verse}

Nolan captura toda la escena del viaje al Hades con envidiable precisión y economía. Los versos que describen el averno son tremendos y las imágenes del film, entreveradas de oscuridad, lo reflejan fielmente. Este episodio, junto al de las sirenas, está a la altura de la historia que cuenta.

Odiseo, después de recibir las instrucciones, se resigna a hacer el viaje y avisa a sus hombres. Pero no van a marcharse de la isla de Circe sin sufrir otra desgracia:

\begin{verse}
Eso les dije, y su aliento y arrojo a mi voz se rendían.\\
Mas ni de allí pude ilesos sacar a mis hombres. Había\\
uno, el más joven de todos, Elpénor, que ni era de hombría\\
para la guerra excesivo ni por su cordura lucía,\\
que por el santo palacio de Circe y sin compañía,\\
fresca buscando y pesado de vino, a dormir se fue arriba.\\
Cuando las voces oyó y el trajín, porque ya se movían\\
sus compañeros, saltó de repente sin ver que debía\\
ir hacia atrás y bajar por la larga escalera: de crisma\\
desde el tejado cayó y su cuello quedó por la espina\\
roto y quebrado, y al Hades se hundió su resuello de vida.
\end{verse}

Una muerte absurda, que conmueve por lo arbitrario, lo casual, lo fuera de lugar. Parece que Homero quiera recordarnos que la parca nos espera detrás de cada esquina. El pobre Elpénor, el más joven de los compañeros, aparece en escena solo para que las furias lo arrojen, de un zarpazo, al Hades.

El canto termina de manera misteriosa, con Circe paseándose invisible por la nave de los aqueos:

\begin{verse}
Cuando a la léveda nave, a la orilla del mar, afligidos,\\
lágrima en flor derramando hacíamos todos camino,\\
Circe nos vino a pasar por delante y al negro navío\\
ató una oveja bien negra con un carnero chiquito, y\\
volvió a pasar sin ser vista de largo: a dios que no quiso\\
¿quién con sus ojos lo vio alguna vez cuando fue y cuando vino?
\end{verse}

¿Quién es Circe y qué pasó exactamente en Eea? No hace falta que repita que Ulises no es un narrador fiable y es bien posible que la historia que le está contando a los feacios incluya omisiones, fantasías o fabulaciones. ¿Por qué la hechicera, después de intentar convertirlo en cerdo como a todos sus hombres, accede tan rápidamente a sus deseos? ¿De verdad está asustada por las amenazas de un mortal? ¿Tiene miedo cuando le abraza las rodillas, o está jugando con él por divertirse? Ciertamente Odiseo no se hace de rogar para subir con ella al lecho, ni le parece tiempo excesivo pasarse un año en la isla. De hecho, son sus hombres los que tienen que suplicarle que sigan viaje.

Quizás hay otra lectura de Circe. En el epílogo de \emph{Abandonando Ítaca}, Candelas Gala escribe:

\begin{quotation}
La elección de seguir la dirección de Circe rompe con las expectativas. Ulises renuncia a la seguridad del hogar en Ítaca, a la lealtad y fidelidad de una esposa, por una situación tan potencialmente sorprendente e inesperada como el cambio del mar. Penélope le servirá un tejido/texto de hilos predecibles cuya linealidad acabará en la muerte; pero los nudos y giros de Circe le llevarán por caminos tortuosos, impredecibles, incluso peligrosos, pero vivos, donde lo inesperado contiene una plenitud vital que implica azar y probabilidad, pero también creatividad permanente.
\end{quotation}

En mi versión del poema, Ulises regresa a Ítaca, como se espera de él, pero es incapaz de olvidar a la hechicera. Muchos años más tarde, desde su todavía ardiente senectud, escribe:

\begin{verse}
Ha pasado mucho tiempo.\\
Intentas invocar su rostro, y lo único que puedes recordar\\
son sus ojos ardientes y su ávida boca.\\
Nunca ha hablado mucho, esa bruja,\\
ni una sola vez dijo\\
que te amaba.
\end{verse}

En realidad, tampoco se lo dijo nunca Penélope. Una cosa es el sexo y otra declarar los sentimientos. Nadie en la \emph{Odisea} habla nunca de amor.

Pero el recuerdo más lacerante para el anciano es la voz de Circe:

\begin{verse}
A menudo, la recuerdas cantando,\\
como aquella vez, cuando cantó en el entierro de uno de tus hombres,\\
era un don nadie, un pobre soldado que una noche bebió demasiado\\
y se subió al tejado, quizás en un vano intento\\
de ver desde más cerca las estrellas, perdió el pie\\
y se mató al caer, lo encontrasteis más tarde,\\
vagando por el Hades, la sombra de una sombra,\\
no más insignificante en la muerte de lo que fue en su vida.

Y sin embargo, él también, una vez, estuvo vivo.
\end{verse}

\begin{verse}
Circe lo vio, ella que no tenía paciencia con los mortales,\\
y aun menos con los hombres,\\
se sintió, sin embargo, conmovida por el hecho ineludible\\
de que el pobre muchacho, tendido en su mesa, con el cuello roto\\
y el asombro en sus ojos, como si en su último momento\\
hubiese captado una verdadera visión de las constelaciones,\\
aquel niño, que yacía para siempre inmóvil, aquella cosa cuya forma\\
no seguiría siendo humana mucho tiempo, ayer mismo,\\
estuvo respirando bajo el Sol.

Y cuando cantaba, su voz era tan pura,\\
subiendo muy arriba, hasta las esferas,\\
cantando en la noche, como si nadie antes de ella\\
hubiese cantado, o llorado.
\end{verse}

Circe, en el recuerdo del viejo Ulises, deja de ser una diosa altiva, distante y no poco malvada, para convertirse en una mujer capaz de llorar por una muerte absurda, por todas las muertes absurdas que abundan en la \emph{Odisea}:

\begin{verse}
Lloró por el joven que yacía muerto a sus pies,\\
y por todos los jóvenes que perecieron en Troya,\\
o que fueron borrados por Cíclopes y Lestrigones,\\
por sirenas y tormentas, cortados, quemados, hundidos,\\
con sus miembros perfectos cruelmente mutilados,\\
y la luz, para siempre robada de sus ojos.

Lloró también por los ancianos, como tu padre,\\
cuya vida ha sido larga,

aunque la vida humana nunca es bastante larga.\\
Los viejos, que van siendo niños de nuevo\\
pues todos los humanos, al envejecer,\\
regresan a su infancia.

Y como ocurre a un niño, el viejo está tan asustado\\
como el muchacho que masacraste aquella noche en Troya,\\
tu espada atravesó su vientre, sus tripas se desparramaron,\\
y trató de sostenerlas con sus manos, y llamó a su madre,\\
no sentiste piedad por él entonces, pero ahora está aquí,\\
contigo (y otros tantos fantasmas),\\
mientras te acuerdas de Circe.
\end{verse}

Ulises, consumido por remordimientos tardíos, se aferra al recuerdo de aquella voz prístina como se aferraba al madero que lo llevó a la isla de Nausícaa.

\begin{verse}
Sí, ella cantó. Sus ojos estaban cerrados, su pecho en alto.\\
En su voz podías escuchar las cuerdas de la lira\\
que tocaba Orfeo. Y la canción también lamentaba a Eurídice,\\
que no pudo ser rescatada del Inframundo,\\
y se lamentó por la esposa de Héctor a la que violaste,\\
y se lamentó por todas esas otras mujeres que fueron tomadas,\\
por carniceros como tú.
\end{verse}

Y es esa voz la que le trae claridad, la que le desvela quién es, quién ha sido siempre:

\begin{verse}
La escuchaste cantar, y su voz era hermosa y, a un tiempo,\\
un dedo señaló tu pecho, preguntando por qué.\\
La escuchaste cantar, preguntándote\\
cómo es que un monstruo como tú merecía tanta hermosura.\\
Un monstruo como tú, que no debería tener corazón,\\
sino una fría piedra bombeando su negra sangre, y sin embargo,\\
abrumado por la pena, apenas podías respirar.

Sí, ella cantó. Su voz era un espejo, donde veías tus crímenes,\\
pero también una oración compasiva,\\
y un ofrecimiento de perdón. Su voz dijo\\
monstruo o no, eres solo un hombre, y todos los hombres,

lo merezcan o no,\\
pueden ser redimidos.
\end{verse}



\end{document}
